\section{Higher determinant products as local-coefficient connecting maps}
\label{sec:136determinant}

\paragraph{Archive and source scope.}
The application is to cote 136, physical pp.~117,119,121, archive
footers 116,118,120, supplied scans 117,119,121.
The external input is the definition of the higher determinant product
in Flynn--Connolly, Moreno--Fernández and Muro,
\emph{Determinant Massey products}, arXiv:2611.25690v1
(19 November 2026), Definition 5.1 and equation (6).
Their Theorem 5.3 proves that its output is closed, and Proposition 5.4
states naturality. We use only their all-degree-one specialization;
the comparison below is proved here.
\url{https://arxiv.org/html/2611.25690v1}

Fix a commutative ring \(k\), \(S=\Gamma_k(T)\), and the admissible
strictly commutative theories of the archive. All auxiliary parameter
modules below are finite free over \(S\). They are not asserted to be
new admissible PD theories.

\paragraph{ARCH-136-TWISTED-COEFFICIENT-COMPARISON.}
The canonical bounded-degree cohomology comparison for an admissible
theory extends to the nilpotent coefficient systems constructed below,
with a weight cutoff independent of their rank and nilpotence length.
We first give the homological lemma that supplies the required uniformity.

\emph{Homological lemma.}
Let \(R\) be any ring, let \(K\) be a nonnegative homological complex of
projective \(R\)-modules, and let \(M\) be an acyclic nonnegative
homological complex. Then
\[
 \operatorname{Ext}^r_{\operatorname{Ch}_{\geq0}(R)}(K,M)=0
                         \quad(r>0).                 \tag{136.DC1}
\]
For a module \(V\), regarded as a complex in degree zero, the same
derived Hom is computed by the ordinary coefficient cochain complex
\(\operatorname{Hom}_R(K,V)\).

\begin{proof}
Put
\[
 Q(K)=S_0(K_0)\oplus\bigoplus_{j\geq1}D_j(K_j)
                  \longrightarrow K .
\]
On an upper disk component the map is the identity; on its lower
component it is the differential of \(K\). The map is degreewise split
onto, its source is projective in the category of nonnegative
complexes, and its kernel \(K'\) has
\(K'_j=K_{j+1}\) and differential \(-d_K\).

Every chain map \(f:K'\to M\) extends to \(Q(K)\to M\).
Indeed choose maps \(h_j:K_j\to M_j\) recursively so that
\[
             d_Mh_{j+1}-h_jd_K=f_j .
\]
Start with \(h_0=0\). At step \(j\), the map
\(f_j+h_jd_K:K_{j+1}\to M_j\) takes values in \(Z_jM\):
the chain-map equation for \(f\), together with the previous step,
makes its differential zero. Acyclicity makes
\(M_{j+1}\to Z_jM\) surjective, and projectivity of \(K_{j+1}\)
provides the lift. This proves the extension assertion, hence
\(\operatorname{Ext}^1(K,M)=0\). Repeat on the shifted tails \(K'\)
and use dimension shifting to get (136.DC1).

Iteration of \(Q\) gives a projective resolution of \(K\).
Applying \(\operatorname{Hom}(-,V[0])\), every positive disk term has
zero Hom, while the degree-zero summand at resolution stage \(r\)
gives \(\operatorname{Hom}_R(K_r,V)\). The resulting differential is
the ordinary coefficient differential, up to the harmless standard
resolution sign. This proves the second assertion.
\end{proof}

Here is the local-coefficient consequence. Let \(P\to X\) be a principal
simplicial \(G\)-cover, let \(R=k[G]\), and put \(K=N(k[P])\).
Each \(K_j\) is free over \(R\): choose one representative in each orbit
of nondegenerate simplices of \(P\).
Let \(W\) be a graded finite free \(S\)-module, with a homogeneous
\(S\)-linear action of \(G\). Suppose all the weights of \(C\) occurring
in the fixed-total-weight complex
\[
             M_C(W)=(C\otimes_S W)^{v,*}              \tag{136.DC2}
\]
are at least a common bound \(N\). If \(N\) exceeds the vertical and
horizontal bounds of \(C\) through form degree \(d+1\), then there is a
natural comparison, an isomorphism through degree \(d\),
\[
 H^j\operatorname{Hom}_{sR}(k[P],M_C(W))
     \ \cong\ H^j(X;\mathscr V_W),\qquad
 \mathscr V_W=P\times_G W^v .                        \tag{136.DC3}
\]
It is simultaneous for \(G\)-equivariant maps of the modules \(W\).

\begin{proof}
Use the projective resolution of \(K\) in the preceding proof,
transported by Dold--Kan to \(sR\), and take its Hom double complex
against \(M_C(W)\).
As an underlying \(k\)-module, each form row of (136.DC2) is a finite
direct sum of rows \(C^{u,q}\), all with \(u\geq N\).
The horizontal condition of \(C\) makes its normalization acyclic.
Equation (136.DC1) consequently makes these rows acyclic for higher
derived sections, so ordinary sections and derived sections have the
same cohomology through degree \(d\).

The coefficient inclusion
\[
                 \underline{W^v}\longrightarrow M_C(W)
\]
is \(G\)-equivariant. In every original simplex, it is a cohomology
isomorphism through degree \(d+1\), by the vertical condition and the
finite free decomposition of \(W\). Hom from the projective resolution
is exact. The vertical-first spectral sequence therefore identifies
derived sections with the coefficient cohomology through degree \(d\).
The preceding homological lemma identifies that coefficient complex
with \(\operatorname{Hom}_{k[G]}(N(k[P]),W^v)\), the ordinary complex
for the associated local system. Both spectral sequences have finite
diagonals in each total degree.

The argument uses the same projective resolution for every \(W\), so
it respects all their maps and the derived-section connecting maps
from short exact sequences. If the ordinary-section sequence is also
degreewise exact, its connecting maps agree under the comparison.
That exactness holds for the twisted sequences used below, where it
is visible before applying the gauge. The cutoff depends on the form range and on the least
occurring weight, not on the number of free summands of \(W\), on \(G\),
or on the simplicial set \(X\). Products or cochains on infinite \(X\)
are never interchanged with a scalar localization.
\end{proof}

\paragraph{ARCH-136-DETERMINANT-LINEARIZATION.}
Let \(a_i,x_i\) be closed forms of cochain degree one, \(1\leq i\leq n\),
\(n\geq3\), with compatible weights
\[
 \operatorname{wt}(a_i)=w_i,\qquad
 \operatorname{wt}(x_i)=v+w_i.                       \tag{136.DC4}
\]
For the moment work in any strictly graded-commutative cochain algebra
in which these elements occur. Define
\[
 \mathcal R=S[t_1,\ldots,t_n]/(t_1^2,\ldots,t_n^2),\quad
 |t_i|_{\rm form}=0,\quad \operatorname{wt}(t_i)=-w_i,
 \quad A=\sum_i a_it_i,\quad D=d+A\cdot .
\]
Since \(dA=0\) and \(A^2=0\), this is a differential on the associated
module complexes. It is a module differential, not an asserted algebra
derivation. Put
\[
 I=(t_1,\ldots,t_n),\quad
 J=S\,t_1\cdots t_n,\quad M=I/J,\quad L=I/I^2 .
\]
Both \(J\) and \(L\) have their ordinary differential \(d\), since
multiplication by \(A\) is zero on them. There is a short exact sequence
of module complexes and a compatible projection
\[
 0\longrightarrow(C\otimes_SJ)^v
 \longrightarrow(C\otimes_SI)^v
 \longrightarrow(C\otimes_SM)^v\longrightarrow0,
 \qquad M\longrightarrow L .                        \tag{136.DC5}
\]

For fixed representatives \(a_i\), the full determinant product set is,
up to one sign depending only on \(n\),
\[
 \delta\left\{z\in H^1((C\otimes_SM)^v,D):
          \pi z=([x_1],\ldots,[x_n])\in
                  \bigoplus_iH^1(C^{v+w_i,*})\right\}
   \subseteq H^2(C^{v+\sum_iw_i,*}).                 \tag{136.DC6}
\]
Here \(\delta\) is the connecting map from (136.DC5).
The set is empty or a single coset of a subgroup, since its description
is the image of a fiber of homomorphisms of abelian groups.

\begin{proof}
For a subset \(S'\) of cardinality \(r\), the sign in the cited
Definition 5.1 is \((-1)^{r+1}\), because all \(a_i\) have degree one.
Choose signs \(s_2=1\) and
\(s_r=(-1)^{r+1}s_{r-1}\) for \(r\geq3\).
Put \(\sigma_{\{i\}}=x_i\) and
\(\sigma_{S'}=s_r\zeta_{S'}\) for \(r\geq2\).
The complete defining equations then become
\[
             d\sigma_{S'}=
                   \sum_{i\in S'}\sigma_{S'\setminus\{i\}}a_i
                    \qquad(2\leq|S'|<n).
\]
This includes the two-element equation because its two odd terms
anticommute. Writing
\[
       \sigma=\sum_{\varnothing\neq S'\subsetneq[n]}
                         \sigma_{S'}t_{S'},
\]
the equations are exactly \(D\sigma=0\) in \(M\).
The top coefficient of \(D\sigma\), after lifting by a zero top entry,
is minus the sum on the right for \(S'=[n]\).
The paper's output is consequently \(-s_n\delta[\sigma]\).

Conversely, a degree-one \(D\)-cocycle representing any class in the
fiber in (136.DC6) has linear coefficients representing the prescribed
classes \(x_i\), and its remaining coefficients are a defining system.
Taking a different cochain representative changes the connecting value
by a coboundary. This proves the exact full-set formula, including the
case when no defining system exists.

For completeness, the formula includes all choices of representatives
of the \(a_i\). If \(a_i'=a_i+du_i\), the integral unit
\[
                         g=\prod_i(1+u_it_i)
\]
satisfies \(g^{-1}dg=\sum_i du_i\,t_i\).
Multiplication by \(g\) identifies the corresponding twisted module
complexes. It is the identity on \(L\) and on \(J\).
Thus it preserves both the prescribed linear cohomology classes and
the output set. Changing representatives of the \(x_i\) is already
allowed in the cohomological fiber in (136.DC6).
No division by two or by an arity factorial is used.
\end{proof}

\paragraph{ARCH-136-DETERMINANT-RIGIDITY.}
For any two admissible strictly commutative theories \(C,D\) over \(k\),
there is one integer \(N=N(C,D)\) such that every full degree-one
determinant product set agrees under their canonical cohomology
comparisons, for every simplicial set \(X\), every finite arity \(n\geq3\),
and every weight choice (136.DC4) satisfying
\[
                        w_i\geq N,\qquad v+w_i\geq N .
                                                               \tag{136.DC7}
\]
The cutoff is independent of arity, of the offset \(v\), and of \(X\).
In particular these higher determinant products cannot distinguish
integral connected components in the archive problem.

\begin{proof}
Choose \(N\) to dominate the uniform vertical and horizontal bounds of
both theories through form degree three, and enlarge it to be positive.
All nonempty subset weights
\[
                      v+\sum_{i\in S'}w_i
\]
are at least \(N\): choose \(i\in S'\) and use
\(v+w_i\geq N\) and the positivity of all remaining \(w_j\).
Thus every component of each of \(I,J,M,L\) lies in the same comparison
range, independently of the number of subsets.

Fix closed representatives \(a_i\) in \(C(X)\).
At a simplex \(\sigma:\Delta[r]\to X\), define the set
\[
 P_i(\sigma)=\{f\in C_r^{w_i,0}:df=a_i|_\sigma\}.
\]
It is nonempty by vertical \(H^1=0\).
Differences of two such primitives are exactly coefficient constants
\(S^{w_i}\), by vertical \(H^0=S^{w_i}\).
Restriction along every simplex map is an equivariant map between
nonempty torsors for the same constant additive group, hence a bijection.
The product of these torsors is therefore a principal simplicial
\(G\)-cover \(P\to X\), where
\[
                         G=\prod_i(S^{w_i},+).
\]

The class of the \(i\)-th torsor is exactly the canonical ordinary
cohomology class of \(a_i\).
One way to verify this identification is the exact sequence of
simplicial modules
\[
 0\longrightarrow\underline{S^{w_i}}\longrightarrow C^{w_i,0}
          \xrightarrow{d}Z^1(C^{w_i,*})\longrightarrow0.
\]
The pullback over the section \(a_i\) is precisely \(P_i\), and the
connecting class is the coefficient-normalized degree-one de Rham
comparison. Equivalently, choose primitives at vertices; their changes
along edges give the usual torsor one-cocycle, and the triangle identity
is its cocycle equation. Changing vertex primitives changes it by a
coboundary. This also shows directly that equal ordinary classes give
isomorphic principal covers, without restrictions on \(X\).

For \(c=(c_i)\in G\), let
\[
                         \rho(c)=\prod_i(1+c_it_i).
\]
It acts on \(I,J,M,L\) by multiplication and preserves total weight.
Since \(t_i^2=0\), \(\rho(c+c')=\rho(c)\rho(c')\).
Thus these are four compatible coefficient representations of the
same group. Their representations on \(J\) and \(L\) are trivial.

On \(P\) there are tautological primitive zero-forms \(f_i\).
The integral unit \(g_f=\prod_i(1+f_it_i)\) satisfies
\(g_f^{-1}dg_f=A\), so
\[
                         d(g_f\omega)=g_fD\omega .
\]
Under a deck translation by \(c\), \(g_f\) changes to \(\rho(c)g_f\).
Consequently multiplication by \(g_f\) identifies, for each of the four
modules \(W\), the twisted complex on \(X\) with the complex of
\(G\)-equivariant ordinary \(C\)-forms on \(P\) with values in \(W\).
The inverse is multiplication by \(g_f^{-1}\) followed by descent.
Finite freeness of \(W\) justifies commuting it with the section limits.

Apply ARCH-136-TWISTED-COEFFICIENT-COMPARISON, with \(d=2\), to these four
representations. It identifies their cohomology groups through degree
two with the cohomology of the associated local systems on \(X\).
All maps in (136.DC5), their connecting map, and the projection to \(L\)
are identified simultaneously. On \(L\) and \(J\), this is exactly the
ordinary canonical cohomology comparison: the representations are
trivial, and the coefficient inclusions and connecting constructions
are the same normalized coefficient comparison.

Now choose in \(D(X)\) representatives of the same prescribed ordinary
classes \(a_i\) and \(x_i\). Their primitive torsors give an isomorphic
principal \(G\)-cover, by the preceding torsor-class calculation.
The coefficient representations depend only on the weights and on
\(S\), so this cover isomorphism identifies the entire four-term
local-coefficient diagram. On its trivial end representations \(L\)
and \(J\), it induces the identity ordinary cohomology maps.
It therefore identifies the fiber specified by the classes \(x_i\)
and its image under the connecting map.
Formula (136.DC6) proves equality of the full determinant sets.

The choice of cover isomorphism is harmless: any other such isomorphism
still commutes with every representation map and induces the identity
on the two trivial end representations. It consequently has the same
effect on the full image of the specified fiber.
This proves the theorem with one cutoff independent of all finite
arities. The argument uses strict commutativity only in the connection
and its integral gauge; it does not infer preservation from an
associative zigzag.
\end{proof}

\paragraph{Exact limitation.}
This closes the all-degree-one higher determinant family in Definition
5.1, with the compatible weights stated above. It neither proves the
paper's conjectured higher matric relation nor implies that associative
formality reflects commutative formality for arbitrary dg algebras.
It constructs no integral strict-PD common source, common target, or
contraction of the archive's localization groupoid.

